\documentclass[letterpaper,11pt]{article}
\usepackage[T1]{fontenc}
\usepackage[utf8]{inputenc}
\usepackage[hidelinks]{hyperref}
\usepackage{xcolor}
\usepackage[margin=1in]{geometry}
\usepackage[default]{sourcesanspro}
\definecolor{ResumeNavy}{HTML}{18324A}
\definecolor{ResumeSlate}{HTML}{5B6573}
\color{black}
\setlength{\parskip}{8pt}
\setlength{\parindent}{0pt}
\hypersetup{colorlinks=true, urlcolor=ResumeNavy}
\begin{document}
\begin{center}
  {\LARGE \bfseries \textcolor{ResumeNavy}{Deva Anand}} \\[3pt]
  {\small
  \href{mailto:devaanand@umass.edu}{devaanand@umass.edu} \enspace\textbar\enspace
  (413) 315-1715 \enspace\textbar\enspace
  \href{https://github.com/Deva-1903}{github.com/Deva-1903} \enspace\textbar\enspace
  \href{https://linkedin.com/in/devaaa/}{linkedin.com/in/devaaa}}
\end{center}
\vspace{6pt}

June 11, 2026

\vspace{2pt}
Cisco \\
Data \& ML Infrastructure Intern Hiring Team

\vspace{6pt}
Dear Hiring Team,

I'm applying for the Data \& ML Infrastructure Intern role. I'm a Computer Science master's student at UMass Amherst, and what drew me to this posting is that it sits exactly where I like to work: the seam between data engineering and machine learning, building the pipelines and infrastructure that let other teams actually use AI rather than the models themselves.

The work I'm proudest of is in that seam. I built a distributed Spark/PySpark ETL and analytics pipeline over a roughly 1.4-billion-row dataset and treated it as a real performance problem, not a homework exercise: I cut end-to-end runtime about 25\% (from 17--19 minutes to 14.3) with broadcast joins, Adaptive Query Execution, shuffle-partition tuning, and schema-aware column-pruned reads, and brought the analytics tail-latency ratio (P99/P50) from 2.14x down to 1.74x with partition pruning, filter pushdown, and adaptive skew-join handling. I localized the shuffle and skew bottlenecks through Spark UI stage timings and wrote up a repeatable optimization playbook. That mindset, profile first, then tune the part that actually hurts, is what I'd bring to keeping data infrastructure fast and reliable at Cisco's scale.

On the ML-infrastructure side, my two years as a backend engineer at Wysa were largely about building the tooling that supports a machine learning team. I architected a full-stack ML training-data annotation platform that the AI team used daily to label data for production models, which cut over 100 hours of manual annotation work a month and improved labeling throughput by 60\%. Separately I owned a multi-tenant backend serving 8+ UK clients and 10,000+ monthly submissions, so scalability, reliability, and fault tolerance aren't abstract to me, they're things I've had to keep working under load and across releases.

I round that out with the ML and systems fundamentals the role asks for: Python and data structures as a daily language, ML frameworks (PyTorch and scikit-learn, with from-scratch autodiff and training loops in my graduate ML coursework), and hands-on Docker and Kubernetes for packaging and deploying workloads across AWS, GCP, and Azure. I'm genuinely excited by Cisco's bet on data and intelligent automation across business groups, and I'd be glad to contribute to the platforms that make it real.

Thank you for considering my application.

\vspace{8pt}
Sincerely, \\
Deva Anand
\end{document}
